% ===== §4 The Symbolic Order and Its Two Fates =====
\section{The Symbolic Order and Its Two Fates}
\label{sec:symbolic}

\noindent
A debt was left unpaid in the diagnosis of the everyday. The third pathology, the decorative symbolisation that photographs the meal and narrates the moment, was named and then deliberately set aside, on the ground that it could not be analysed until the theory of the symbolic order was in hand (\xref{sub:diag-pathologies}). This section pays that debt, and in paying it corrects a crudeness that the word \emph{pathology} invited. It would be easy, and it would be wrong, to treat the symbolic intervention in perception as such a pollution, a fall from some purer immediate seeing. The symbolic order is the condition of shared perception, not its enemy; without symbolisation there is no perception two people could hold in common, only the silted-up interiority of a private imaginary. The symbol is therefore not the disease. What the diagnosis glimpsed as a pathology is one of two fates that the symbolic intervention may meet, and the work of this section is to distinguish the fate in which symbolisation generates from the fate in which it appropriates, and to give the criterion that tells them apart.

\subsection{Symbolisation as the condition of shared perception}
\label{sub:sym-condition}

Perception, left in the register of the imaginary alone, is unshareable. The account of the three registers this series has used throughout locates in the imaginary the domain of the perceiver's own images, vivid and immediate and idiosyncratic, and it is a standing feature of that register that its contents cannot as such be given to another; my image of this light is mine, and remains mine, until it is articulated in a system that another can enter. That system is the symbolic order, and its defining property, since Saussure, is that it signifies by difference rather than by resemblance: a term means through its place in a structure of contrasts, not through any likeness to what it marks \citep{saussure1959}. It is precisely this differential structure that makes the symbolic shareable where the imaginary is not, for a structure of differences can be occupied by more than one subject, whereas an image cannot be inhabited by anyone but the one who has it. When a perception is symbolised, it is lifted out of the private immediacy in which it could only be undergone and given a form in which it can be transmitted, revisited, and held in common. The intervention of the symbolic in perception is thus the very operation by which a perception becomes something that a relation, and not only a solitary perceiver, can have. To wish the symbol away in the name of an unmediated seeing is to wish away the possibility of a shared aesthetic life, which is the possibility this paper is about. The symbolic is the condition of the relational, and the relational is where, on this paper's central claim, beauty is generated. The question is therefore never whether perception should be symbolised. It is always what becomes of a perception when it is.

\subsection{The generative fate}
\label{sub:sym-generative}

A symbolisation may generate. It does so when the form it gives a perception opens that perception to renewed generation rather than closing it into a token, when the symbol returns the perceiver, and the one to whom the perception is transmitted, to the generative encounter rather than substituting for it. The mark of the generative fate, in the vocabulary this series has built, is that the symbolisation is organised around the empty place and does not settle it. A perception symbolised generatively keeps the incommensurable core of what was perceived unfilled, presenting a form that points toward the real value at its centre without claiming to have captured it, and so remains an invitation to perceive again rather than a receipt for having perceived. The love letter that says the day was beautiful in a way that sends its reader back to the day, the shared word between two people that names a recurring quality of their evenings and thereby deepens their attention to it, the poem that does not exhaust the light it speaks of but teaches the eye to return to it: these are symbolisations whose fate is generative, because each returns its participants to the generative relation and leaves the real value at its heart uncommandeered. Such symbolisation is, moreover, the very medium of relational aesthetic life, since it is only through symbols of this kind that the perception generated between two people can be sedimented into a form they can revisit together. Far from being a fall, the generative symbol is how a shared aesthetics accumulates. The paper will return to this in its account of the benevolent sign, where the symbolic arrangement of a shared life becomes one of the practical models of everyday aesthetic generation (\xref{sec:typology}).

\subsection{The appropriative fate}
\label{sub:sym-appropriative}

A symbolisation may instead appropriate, and this is the fate the diagnosis glimpsed. It does so when the symbol is substituted for the perception rather than opened onto it, when possessing the token of having-seen takes the place of the seeing. The photograph taken so that the meal need not be tasted, the moment narrated so that it need not be undergone, the ordinary arranged into a legible sign of itself for display and archived in place of being lived: in each the symbol closes over the perception it was meant to carry, settling the real value into a commensurable and exchangeable mark. The appropriative symbol belongs, in the value theory of this series, to the register that is most readily stolen, the symbolic value that can be counted, compared, accumulated, and displayed. Its social form is the one this paper will later confront under its own name, the spectacle, in which perception is handed over to the sign and the generative encounter is replaced by the consumption of images; that confrontation belongs to the chapter on justice, where the political economy of the appropriative symbol can be given its full weight (\xref{sec:distribution}). What must be fixed here is only the structure of the appropriative fate as such. It is not that a symbol has been used, for the generative fate also uses symbols. It is that the symbol has been settled, closed over the empty place it should have kept open, converted from an invitation to perceive into a possession that dispenses with perceiving. The appropriative symbol does not transmit a perception; it retires one.

\subsection{The criterion}
\label{sub:sym-criterion}

The two fates must be told apart by a criterion that does not appeal to whether a symbol was used, since both use symbols, nor to any inner sincerity of the one who symbolises, since the difference is structural and not a matter of avowed intention. This series has such a criterion already, and the present distinction is its instance in the register of perception. A relational process carries a good or a bad character according to the sign of the holonomy of its value cycle: the good cycle returns to its root with a genuine surplus generated around the empty place, an endogenous increment produced by the internal tension of the relation and possessed by no one, which the formal companion of this series renders as a positive geometric phase \citep{berry1984,paperix}; the bad cycle either generates no surplus across its return or generates one only by external extraction, closing the loop and settling the real value it should have circled. The generative symbolisation is the positive holonomy of this cycle in the domain of shared perception. It returns the perceivers to the encounter with a surplus of attention and shared meaning that neither owns and that was produced by their joint perceiving, and it leaves the incommensurable centre unsettled, so that the cycle can turn again. The appropriative symbolisation is the vanishing or the negativity of that holonomy. It closes the loop, settles the centre into an exchangeable token, and returns the perceivers not to the encounter but to the possession of its sign, from which no further generation issues. The criterion is therefore neither the presence of the symbol nor the psychology of the symboliser but the sign of the holonomy: whether the symbolisation, taken through its full cycle, returns a surplus to the shared perception around an unsettled centre, or retires the perception into a settled and possessable mark. With this the debt of the diagnosis is paid. The decorative symbolisation named among the pathologies of the everyday is not the intervention of the symbol, which is the condition of any shared aesthetics, but the appropriative fate of that intervention, the symbolisation whose holonomy has gone to zero or turned negative; and the cure for it is not the abolition of the symbol but the cultivation of the other fate, the generative symbolisation whose possibility this section has established and whose practice the later parts of the paper will develop.
